\documentclass[../main.tex]{subfiles}
\begin{document}

\section{Background}

\subsection{PSBD as its authors state it}

Li et al. \citep{li2025psbd} define a pair of quantities. Prediction Shift is the event that the class a model predicts changes once dropout is enabled at inference, and the shift ratio $\sigma$ is its frequency over a dataset and over $k$ dropout passes. Prediction Shift Uncertainty is the drop in the confidence of the no-dropout prediction, averaged over the $k$ passes, and a low value marks a poisoned input. The paper's threshold is the \HeadlineQuantilePercent{} quantile of PSU on a clean validation set, read as the tolerable loss rate on clean data, and its dropout rate is chosen by an adaptive rule, the rate at which the clean validation shift ratio approaches a high value, \AdaptiveShiftTarget{} in their experiments. Dropout is placed after each residual connection in the residual basic block, before the activation, and inference uses $k = \ForwardPasses$ passes.

The paper's explanation is what this project set out to test on a transformer, and it is worth quoting rather than paraphrasing. The abstract states that ``poisoned models' predictions on clean data often shift away from true labels towards certain other labels with dropout applied during inference, while backdoor samples exhibit less PS'' and that ``we hypothesize PS results from the neuron bias effect, making neurons favor features of certain classes''. The method section makes the target-class claim explicit: among the samples experiencing Prediction Shift, ``almost all clean data shifts to the target class $y_t$'', and the appendix generalises it, ``in all attack scenarios, clean data exhibits a bias towards the target class''. The mechanism offered is that ``under dropout conditions, many key distinguishing features in clean data are discarded. Consequently, the model relies more heavily on the neuron bias established during training, leading it to classify clean data to the label associated with this bias'', while backdoor data has ``a more stable and pronounced neuron bias'' and so survives the same removal. The supporting evidence is a figure of the final convolutional feature maps of a clean image and its triggered version becoming ``almost identical with dropout''. The paper also records the counterexample on Tiny ImageNet, where ``the shift classes of the clean training data and the clean validation data exhibit a predominant inclination towards a certain class rather than the target class'', and attributes it to inadequate generalisation. Every statement above is a claim about ConvNets, since the paper's experiments never include a Vision Transformer, and its own future-work paragraph names ViT as the open case.

\subsection{Threat model}

We keep the paper's post-hoc setting and change a single thing about the data. The defender holds a trained checkpoint and \HeldoutSize{} clean images with no labels required, may run the model as often as the method needs, and must decide per input whether it carries a trigger, without retraining and without any triggered example. The attacker poisoned the training set at a rate the defender does not know, with a trigger and a label mapping the defender does not know, and the model was trained by the defender's own ordinary recipe. PSBD scores the poisoned training set itself. We score a held-out test pool, the same images clean and triggered, so our numbers answer whether PSBD flags a triggered input at test time and are not directly comparable to the paper's tables, a difference stated once here rather than hedged at every table.

\subsection{The architectures}

ViT-B/16 \citep{dosovitskiy2021an} cuts the image into \VitPatches{} patches, prepends a class token, and runs \VitBlocks{} blocks of width \VitWidth{} with \VitHeads{} heads each, so the classifier reads a single token out of \VitTokens{} at the end. Each block is a read-write cycle on the residual stream, a LayerNorm into attention and an add, then a LayerNorm into the MLP and an add. Swin-S \citep{liu2021swin} has \SwinBlocks{} blocks with windowed attention and no class token, reads the token mean at the end, and shares every position name with ViT through a single registry, so a placement is written once and runs on both.

\end{document}
